% !TeX root = ../../../main.tex
% !TeX encoding = UTF-8

% 第一册第一章第二节：在此编写习题解答。
% 请先在 main.tex 中填写本节的实际标题。
% 以下为写作模板，使用时去掉相应行开头的百分号。

% \begin{exercise}[1]
% \label{ex:b1-c1-s2-q1}
% 在此写题意。
% \end{exercise}
% \begin{solution}
% 在此写解答。
% \end{solution}
\begin{exercise}
\label{ex:b1-c1-s2-q1}
证明集合关系:\\
\begin{align*}
&(A\subset C)\land (B\subset C)\Leftrightarrow ((A\cup B)\subset C)\\
&(C\subset A)\land (C\subset B)\Leftrightarrow (C\subset (A\cap B))\\
&C_M(C_M A)=A\\
&(A\subset C_M B)\Leftrightarrow (B\subset C_M A)\\
&(A\subset B)\Leftrightarrow (C_M A\supset C_M B)
\end{align*}
\end{exercise}

\begin{solution}
\label{sol:b1-c1-s2-q1}
.\\
$(A\subset C)\land (B\subset C)\Leftrightarrow ((A\cup B)\subset C)$:
\begin{align*}
 (A\subset C)\land (B\subset C)&\Leftrightarrow (\forall x\in A, x\in C)\land (\forall x\in B,x\in C)\\
&\Leftrightarrow \forall (x\in A)\lor (x\in B),x\in C\\
&\Leftrightarrow \forall x\in (A\cup B),x\in C\\
&\Leftrightarrow ((A\cup B)\subset C)
\end{align*}
$(C\subset A)\land (C\subset B)\Leftrightarrow (C\subset (A\cap B))$:
\begin{align*}
(C\subset A)\land (C\subset B)&\Leftrightarrow \forall x\in C,x\in A\land x\in B\\
&\Leftrightarrow \forall x\in C,x\in A\cap B\\
&\Leftrightarrow (C\subset (A\cap B))
\end{align*}
$C_M(C_M A)=A$:
\begin{align*}
(C_M(C_M A))&=\{x| x\in M\land x\notin C_M A\}\\
&=\{x| x\in M\land \lnot (x\in C_M A) \}\\
&=\{x| x\in M\land \lnot(x\in M\land x\notin A)\}\\
&=\{x| x\in M \land (\lnot(x\in M)\lor (x\in A))\}\\
&=\{x| x\in M\land \lnot (x\in M)\lor x\in M\land x\in A \}\\
&=\{x| x\in M\land x\in A\}\\
&=A
\end{align*}
$(A\subset C_M B)\Leftrightarrow (B\subset C_M A)$:
\begin{align*}
(A\subset C_M B)&\Leftrightarrow \forall x\in A, x\in C_M B\\
&\Leftrightarrow \forall x\notin C_M B,x\notin A\\
&\Leftrightarrow \forall x\in B,x\in C_M A\\
&\Leftrightarrow (B\subset C_M A)
\end{align*}
$(A\subset B)\Leftrightarrow (C_M A\supset C_M B)$:
\begin{align*}
A\subset B&\Leftrightarrow \forall x\in A,x\in B\\
&\Leftrightarrow \forall x\notin B,x\notin A\\
&\Leftrightarrow \forall x\in C_M B,x\in C_M A\\
&\Leftrightarrow C_M A\supset C_M B
\end{align*}
鉴于本书并非离散数学教材，因此相关命题证明并未完全按照命题语言严谨表述。
\end{solution}

\begin{exercise}
证明集合关系:
\begin{align*}
    &A\cup (B\cup C)=(A\cup B)\cup C\\
    &A\cap (B\cap C)=(A\cap B)\cap C\\
    &A\cap (B\cup C)=(A\cap B)\cup (A\cap C)\\
    &A\cup (B\cap C)=(A\cup B)\cap (A\cup C)
\end{align*}
\end{exercise}
\begin{solution}.\\
$A\cup (B\cup C)=(A\cup B)\cup C$:
\begin{align*}
    x \in A\cup (B\cup C)&\Leftrightarrow x\in A\lor x\in (B\cup C)\\
    &\Leftrightarrow x\in A\lor (x\in B\lor x\in C)\\
    &\Leftrightarrow  (x\in A\lor x\in B)\lor x\in C\\
    &\Leftrightarrow (x\in (A\cup B))\lor x\in C\\
    &\Leftrightarrow x\in (A\cup B)\cup C
\end{align*}
$A\cap (B\cap C)=(A\cap B)\cap C$:
\begin{align*}
    x \in A\cap (B\cap C)&\Leftrightarrow x\in A\land x\in (B\cap C)\\
    &\Leftrightarrow x\in A\land (x\in B\land x\in C)\\
    &\Leftrightarrow  (x\in A\land x\in B)\land x\in C\\
    &\Leftrightarrow (x\in (A\cap B))\land x\in C\\
    &\Leftrightarrow x\in (A\cap B)\cap C
\end{align*}
$A\cap (B\cup C)=(A\cap B)\cup (A\cap C)$:
\begin{align*}
    x\in A\cap (B\cup C)&\Leftrightarrow x\in A \land x\in (B\cup C)\\
    &\Leftrightarrow x\in A\land (x\in B\lor x\in C)\\
    &\Leftrightarrow (x\in A\land x\in B)\lor (x\in A\land x\in C)\\
    &\Leftrightarrow x\in (A\cap B)\lor x\in (A\cap C)\\
    &\Leftrightarrow x\in (A\cap B)\cup (A\cap C)
\end{align*}
$A\cup (B\cap C)=(A\cup B)\cap (A\cup C)$:
\begin{align*}
    x\in A\cup (B\cap C)&\Leftrightarrow x\in A \lor x\in (B\cap C)\\
    &\Leftrightarrow x\in A\lor (x\in B\land x\in C)\\
    &\Leftrightarrow (x\in A\lor x\in B)\land (x\in A\lor x\in C)\\
    &\Leftrightarrow x\in (A\cup B)\land x\in (A\cup C)\\
    &\Leftrightarrow x\in (A\cup B)\cap (A\cup C)
\end{align*}
\end{solution}

\begin{exercise}
证明以下命题:
\begin{align*}
    &C_M(A\cup B)=C_M A\cap C_M B\\
    &C_M(A\cap B)=C_M A\cup C_M B
\end{align*}
\end{exercise}
\begin{solution}.\\
$C_M(A\cup B)=C_M A\cap C_M B$:
\begin{align*}
    x\in C_M(A\cup B)&\Leftrightarrow \lnot (x\in A\cup B)\\
    &\Leftrightarrow \lnot (x\in A\lor x\in B)\\
    &\Leftrightarrow \lnot (x\in A)\land \lnot (x\in B)\\
    &\Leftrightarrow x\in C_M A\land x\in C_M B\\
    &\Leftrightarrow x\in (C_M A\cap C_M B)
\end{align*}
$C_M(A\cap B)=C_M A\cup C_M B$:
\begin{align*}
    x\in C_M(A\cap B)&\Leftrightarrow \lnot (x\in A\cap B)\\
    &\Leftrightarrow \lnot (x\in A\land x\in B)\\
    &\Leftrightarrow \lnot (x\in A)\lor \lnot (x\in B)\\
    &\Leftrightarrow x\in C_M A\lor x\in C_M B\\
    &\Leftrightarrow x\in (C_M A\cup C_M B)
\end{align*}
\end{solution}

\begin{exercise}
    给出以下集合的笛卡尔积的几何解释:
    \begin{align*}
        & a)\text{两个线段}\\
        & b)\text{两个直线}\\
        & c)\text{直线与圆周}\\
        & d)\text{直线与圆面}\\
        & e)\text{两圆周}\\
        & f)\text{圆周与圆面}
    \end{align*}
\end{exercise}
\begin{solution}.\\
    \begin{align*}
        & a)\text{矩形}\\
        & b)\text{平面}\\
        & c)\text{圆柱面}\\
        & d)\text{圆柱体}\\
        & e)\text{圆环面}\\
        & f)\text{圆环体}
    \end{align*}
\end{solution}
\begin{exercise}
    给出对应的笛卡尔平方对角线的几何解释
\end{exercise}
\begin{solution}.\\
    \begin{align*}
        &a)\text{对角线}\\
        &b)\text{平面上y=x的直线}\\
        &e)\text{圆周}
    \end{align*}
\end{solution}


\begin{exercise}
    证明以下命题:
    \begin{align*}
        &a) (X\times Y=\varnothing )\Leftrightarrow (X = \varnothing)\lor (Y=\varnothing)\\
        &\text{在}X\times Y\neq \varnothing, A\times B\neq \varnothing \text{的前提下，证明下列命题}\\
        &b) (A\times B\subset X\times Y)\Leftrightarrow (A\subset X)\land (B\subset Y)\\
        &c) (X\times Y)\cup (Z\times Y)=(X\cup Z)\times Y\\
        &d) (X\times Y)\cap (X^\prime \times Y^\prime)=(X\cap X^\prime)\times (Y\cap Y^\prime)  
    \end{align*}
\end{exercise}
\begin{solution}.\\
$(X\times Y=\varnothing )\Leftrightarrow (X = \varnothing)\lor (Y=\varnothing)$:
\begin{align*}
(X\times Y=\varnothing )&\iff \{(x,y)|x\in X,y\in Y\}=\varnothing\\
&\iff X=\varnothing \lor Y=\varnothing
\end{align*}
$(A\times B\subset X\times Y)\Leftrightarrow (A\subset X)\land (B\subset Y)$:
\begin{align*}
    (A\times B\subset X\times Y)&\Leftrightarrow \forall (x,y),x\in A,y\in B,(x,y)\in X\times Y\\
    &\Leftrightarrow \forall (x,y),x\in A,y\in B, \exists (c,d),c\in X,d\in Y,(x,y)=(c,d)\\
    &\Leftrightarrow \forall x\in A,x\in X,\forall y\in B,y\in Y\\
    &\Leftrightarrow (A\subset X)\land (B\subset Y)
\end{align*}
$(X\times Y)\cup (Z\times Y)=(X\cup Z)\times Y$:
\begin{align*}
    (x,y)\in ((X\times Y)\cup (Z\times Y))&\iff (x,y)\in X\times Y\lor (x,y)\in Z\times Y\\
    &\iff (x\in X\land y\in Y)\lor(x\in Z\land y\in Y)\\
    &\iff (x\in X\lor x\in Z)\land y\in Y\\
    &\iff (x\in X\cup Z)\land y\in Y\\
    &\iff (x,y) \in (X\cup Z)\times Y
\end{align*}
$(X\times Y)\cap (X^\prime \times Y^\prime)=(X\cap X^\prime)\times (Y\cap Y^\prime)$:
\begin{align*}
    (x,y)\in (X\times Y)\cap (X^\prime \times Y^\prime)&\iff (x,y)\in (X\times Y)\land (x,y)\in (X^\prime \times Y^\prime)\\
    &\iff (x\in X\land y\in Y)\land (x\in X^\prime\land y\in Y^\prime)\\
    &\iff (x\in X)\land (y\in Y)\land (x\in X^\prime)\land (y\in Y^\prime)\\
    &\iff (x\in X\cap X^\prime)\land (y\in Y\cap Y^\prime)\\
    &\iff (x,y)\in (X\cap X^\prime)\times (Y\cap Y^\prime)
\end{align*}
\end{solution}
\begin{exercise}
    对比习题3中的关系以及与$\S 1$习题二中的关系a),b),建立起命题逻辑运算$\lnot,\land,\lor $与集合运算$C,\cap,\cup $的关系
\end{exercise}
\begin{solution}.\\
    可以对应到: 逻辑否即为取补集，逻辑或即为取并集，逻辑与即为取交集
\end{solution}
